% !TeX root = tesis.tex

Una prueba de conocimiento cero (\ZeroKnowledgeProof) o prueba \ZK es un protocolo entre dos partes, en donde una quiere convencer a otra de que posee un valor, sin revelar nada sobre el conocimiento de este valor más allá de la veracidad de la información. Estas partes se conocen como \textbf{Prover} $\Prover$ y \textbf{Verifier} $\Verifier$. Podemos pensar a estas partes como personas ejecutando un algoritmo, o como algoritmos en sí mismos (usaremos estos conceptos de forma indistinta). Esta idea fue introducida inicialmente por Goldwasser, Micali y Rackoff en \cite{DBLP:journals/siamcomp/GoldwasserMR89}.

\begin{definition}[Protocolo de pruebas \ZK]
	Sea $\Relation \subseteq \InstanceSpace \times \WitnessSpace$ una relación, donde $\InstanceSpace$ es el conjunto de \textit{public inputs} y $\WitnessSpace$ el conjunto de \textit{witnesses}. Asociado a $\Relation$, definimos el lenguaje
	$$\Language_\Relation := \{x \in \InstanceSpace : \exists w \in \WitnessSpace,\ (x,w) \in \Relation\}.$$
	Un protocolo de pruebas \ZK para $\Relation$ está dado por un prover $\Prover$ y un verifier $\Verifier$ tales que, sobre una instancia pública $x \in \InstanceSpace$ y un witness $w \in \WitnessSpace$, generan una interacción verificable por el verifier. Si el protocolo es interactivo, el objeto matemático relevante es el \textit{transcript} completo de esa interacción; si es no interactivo, dicho objeto se reduce a una prueba $\pi \in \ProofSpace$.

	Informalmente, el protocolo debe satisfacer las siguientes propiedades:
	\begin{enumerate}[noitemsep]
		\item Completitud: si $(x,w) \in \Relation$ y ambas partes siguen el protocolo, entonces $\Verifier$ acepta con probabilidad alta.
		\item Solidez: si la afirmación pública es falsa, es decir, si $x \notin \Language_\Relation$, entonces un prover deshonesto no puede hacer que $\Verifier$ acepte salvo con probabilidad negligible.
		\item Conocimiento cero: el verifier no aprende nada esencial sobre $w$ más allá de la validez de la afirmación demostrada.
	\end{enumerate}

	Diremos que el protocolo es \textit{no interactivo} si toda la comunicación de $\Prover$ a $\Verifier$ consiste en un único mensaje $\pi \in \ProofSpace$; en ese caso llamaremos prueba a dicho mensaje. En caso contrario, diremos que es \textit{interactivo} y hablaremos del \textit{transcript} completo de la ejecución.
\end{definition}

Vamos a ganar intuición sobre esta definición, apoyados también en el protocolo de ejemplo de la imagen \ref{fig:interactive-zk-protocol}. Supongamos que tenemos un valor secreto $w \in \Strings{*}$ y calculamos $x = \Hash(w) \in \Strings{n}$ para una función de hash criptográfica $\Hash$, de modo que $x$ es interesante (por ejemplo, empieza con muchos ceros) y por lo cual nos darían un premio al haber hallado $w$. Si alguien nos pide ver $w$ para verificar la igualdad $\Hash(w) = x$, podemos no compartirlo, pero demostrar que lo conocemos (por esto se habla de prueba de conocimiento). En este contexto:

\begin{itemize}[noitemsep]
	\item La relación $\Relation$ está dada por los pares $(x,w)$ tales que $\Hash(w) = x$. Notar que pueden existir múltiples witnesses válidos para el mismo $x$. Se puede describir de forma concisa como $\Relation = \{(\Hash(w), w): w \in \WitnessSpace\}$.
	\item El conjunto $\InstanceSpace = \Strings{n}$ sería el de \textit{public inputs}, ya que ambos conocemos $x$.
	\item El conjunto $\WitnessSpace = \Strings{*}$ sería el \textit{witness}, ya que solamente nosotros conocemos $w$.
	\item Hay un conjunto de pruebas $\ProofSpace$ que describe todas las posibles combinaciones de datos que le enviamos al verifier. Veremos cómo son estos conjuntos cuando vayamos a protocolos de pruebas ZK concretos.
	\item Pensemos que si un $w$ tal que $(x,w) \in \Relation$ fuera fácil de calcular, el protocolo de pruebas ZK no tendría sentido.
\end{itemize}

En la imagen \ref{fig:interactive-zk-protocol}, el protocolo de ejemplo muestra una interacción de tres pasos entre $\Prover$ y $\Verifier$. En este caso, el objeto relevante no es una prueba aislada sino el \textit{transcript} completo de la interacción, formado por todos los mensajes intercambiados.

\begin{figure}[!t]
	\centering
	\begin{tikzpicture}[font=\small]
		
		\tikzstyle{block} = [rectangle, draw, text centered, minimum height=1cm, minimum width=3.5cm]
		\tikzstyle{msg}   = [->, thick, dashed]
		
		% Nodes (algorithms)
		\node (P) at (0,0) [block] {Prover ($\Prover$), $\InstanceSpace, \WitnessSpace$};
		\node (V) at (7,0) [block] {Verifier ($\Verifier$), $\InstanceSpace$};
		
		% Vertical lifelines
		\draw (P.south) ++ (-0.1, 0) -- ++(-0,-6);
		\draw (V.south) ++ (0.1, 0)-- ++(0,-6);
		
		% Messages (Verifier -> Prover)
		\draw[msg] (7,-1) -- (0,-1) node[pos=0.15, above] {$r_1$};
		\draw[msg] (7,-3) -- (0,-3) node[pos=0.15, above] {$r_2$};
		\draw[msg] (7,-5) -- (0,-5) node[pos=0.15, above] {$r_3$};
		
		% Messages (Prover -> Verifier)
		\draw[msg] (0,-2) -- (7,-2) node[pos=0.15, above] {$\pi_1$};
		\draw[msg] (0,-4) -- (7,-4) node[pos=0.15, above] {$\pi_2$};
		\draw[msg] (0,-6) -- (7,-6) node[pos=0.15, above] {$\pi_3$};
		
		% Output
		\node at (0,-7) {$\textnormal{Transcript}(x,w) = (r_1,\pi_1,r_2,\pi_2,r_3,\pi_3)$};
		\node at (7,-7) {$\Verifier(x,\textnormal{Transcript}) \in \{0,1\}$};
		
	\end{tikzpicture}
	
\caption{Prototipo de un protocolo \ZK interactivo entre $\Prover$ y $\Verifier$ sobre una instancia pública $x$ y un witness privado $w$. Las flechas $r_i$ representan desafíos del verifier, las flechas $\pi_i$ representan respuestas del prover, y el objeto verificable final es el \textit{transcript} completo $(r_1,\pi_1,r_2,\pi_2,r_3,\pi_3)$.}
\label{fig:interactive-zk-protocol}
\end{figure}

Por otra parte, podemos desglosar las tres propiedades que debe cumplir un protocolo de pruebas \ZK.

\begin{itemize}[noitemsep]
	\item Si $\Prover$ ejecuta su algoritmo con un $w$ tal que $(x,w) \in \Relation$, porque efectivamente \textbf{conoce} un $w$ que cumple esto, entonces seguramente podrá convencer a $\Verifier$ de que lo conoce.
		\item Si la afirmación pública es falsa, es decir, si $x \notin \Language_\Relation$, entonces un prover deshonesto $\Prover^*$ probablemente no podrá convencer a $\Verifier$. La probabilidad de engaño no es nula, pero es negligible. Qué tan baja queremos que sea esta probabilidad depende del nivel de bits de seguridad que exigimos a nuestro sistema criptográfico. Observamos que la propiedad de solidez, a diferencia de las otras, vale para \textbf{cualquier} prover.
		\item En esta introducción estamos trabajando con la noción de \textit{honest-verifier zero knowledge}. Esto significa que $\Verifier$ no aprende ninguna información esencial sobre $w$ a partir del \textit{transcript} del protocolo, más allá de la validez de la afirmación, siempre que siga honestamente el protocolo. La propiedad se enuncia a partir de un simulador que no tiene acceso a un witness válido, pero que puede generar una vista indistinguible de una ejecución real para ese verifier honesto. Intuitivamente, la existencia de tal simulador implica que todo lo que ve el \textit{verifier} durante el protocolo podría haberse generado sin conocer el \textit{witness}. Por lo tanto, la interacción no revela información adicional sobre $w$. La clave es la \textbf{indistinguibilidad computacional}. Es crucial que el simulador sea eficiente, pues el tiempo también podría ser una forma de distinguir las interacciones.
		\item Un detalle adicional es que, si bien la propiedad de solidez trata el caso donde un prover puede ser deshonesto, acá no estamos considerando verifiers arbitrarios o maliciosos. La noción usada en esta sección es, precisamente, la de \textit{honest-verifier zero knowledge}.
\end{itemize}

Desde el punto de vista de la teoría de la complejidad, los protocolos de pruebas \ZK interactivos se estudian dentro del marco más general de las pruebas interactivas, correspondiente a la clase \textit{Interactive Proofs} ($\IP$). No todo protocolo en $\IP$ es de conocimiento cero, pero todo protocolo \ZK interactivo es, en particular, una prueba interactiva. Informalmente, una clase de lenguajes $L$ pertenece a $\IP$ si existe un protocolo interactivo $(\Prover, \Verifier)$ tal que para toda instancia $x \in \InstanceSpace$, el verifier $\Verifier$, que se ejecuta en tiempo polinomial, acepta con probabilidad 1 cuando $x \in L$ y rechaza con probabilidad al menos $1 - \epsilon$ cuando $x \notin L$. El prover $\Prover$ no está restringido computacionalmente, y la interacción puede contar con un número polinomial de rondas. Podemos pensar al lenguaje $L$ como inducido por $\Relation \subseteq \InstanceSpace \times \WitnessSpace$, donde $x \in L$ si y sólo si existe $w \in \WitnessSpace$ tal que $(x,w) \in \Relation$.

Un resultado central de la teoría de la complejidad es que la potencia expresiva de los protocolos interactivos coincide exactamente con la de la clase $\PSPACE$. Es decir, todo lenguaje que puede decidirse utilizando una cantidad polinomial de espacio admite un protocolo de prueba interactivo con un verifier probabilístico en tiempo polinomial. En esta tesis, solamente nos interesa hablar del mismo como indicador de la enorme expresividad de los protocolos interactivos.

\begin{theorem}
	$\IP = \PSPACE$.
\end{theorem}

Este teorema fue demostrado inicialmente por Shamir en \cite{Shamir1992}, lo cual implica que problemas como la equivalencia de circuitos booleanos o la validez de fórmulas cuantificadas (\QBF) admiten protocolos interactivos con verificación eficiente. Para lectores experimentados en complejidad computacional, en \cite{arora2006computational} hay una exposición moderna y rigurosa del resultado.

La igualdad tiene implicancias relevantes incluso al restringirse a clases de complejidad más bajas, como $\NP$. En efecto, todo lenguaje en $\NP$ admite trivialmente un protocolo interactivo eficiente: dado un lenguaje inducido por una relación $R \subseteq \InstanceSpace \times \WitnessSpace$, el prover puede convencer al verifier de que $x \in L$ demostrando el conocimiento de un witness $w$ tal que $(x,w) \in \Relation$. Esto permite existencia de protocolos interactivos para problemas clásicos de $\NP$ como la satisfacibilidad booleana, el isomorfismo de grafos, la 3-colorabilidad de un grafo, la existencia de un ciclo hamiltoniano, la existencia de $k$-cliques, entre muchos otros.

Por ahora, consideramos protocolos de pruebas \ZK en su forma más general, y vimos que estos protocolos capturan una clase extremadamente amplia de problemas. Sin embargo, el hecho de que sean interactivos y puedan tener un costo comunicacional elevado los vuelven poco prácticos en muchos escenarios. En particular, los protocolos de pruebas \ZK son usados en sistemas distribuidos y blockchains, donde requieren pruebas que puedan ser verificadas de manera no interactiva, con complejidad de verificación muy baja (no solamente polinomial), y tamaño de prueba reducido.

Más aún, podemos pensar a los protocolos \ZK como una herramienta general para la verificación de cómputo. Dado un algoritmo eficiente y una entrada secreta, el prover puede convencer al verifier de que el algoritmo fue calculado correctamente, sin revelar información adicional sobre $w$ ni sobre el proceso interno de cómputo. Si el tiempo de verificación es bajo, y el tamaño de prueba es reducido, podemos incluso obviar el hecho de que no estamos revelando información sobre $w$ y hablar de protocolos \ZK que permiten verificar cómputos de gran complejidad en sistemas con muy escasos recursos (como una blockchain). También se puede hablar de \textit{cómputo delegado}.

\begin{definition}[Esquema de cómputo verificable (\ZK)]
	Sea $f: \InstanceSpace \times \WitnessSpace \to \mathcal{Y}$ una función computable en tiempo polinomial. Un esquema de cómputo verificable para $f$ es un protocolo $(\Prover, \Verifier)$ asociado a la relación $\Relation_f \subseteq \InstanceSpace^* \times \WitnessSpace$, donde el conjunto de \textit{public inputs} es $\InstanceSpace^* = \InstanceSpace \times \mathcal{Y}$ y el conjunto de \textit{witness} es $\WitnessSpace$, y que está definida por:
	$$(x, y, w) \in \Relation_f \iff f(x,w) = y.$$
	
	El protocolo $(\Prover, \Verifier)$ debe satisfacer:
	\begin{enumerate}[noitemsep]
		\item Completitud: Para todo $x \in \InstanceSpace$ y todo $w \in \WitnessSpace$, si definimos $y := f(x,w)$ y $\pi = \Prover(x,w)$, entonces vale $\Prob[\Verifier(x,y, \pi) = 1] = 1$.
		\item Solidez: Para todo prover probabilístico en tiempo polinomial posiblemente deshonesto $\Prover^\star$ que, dada una entrada pública $(x,y)$, produce una prueba $\pi = \Prover^\star(x,y)$, vale $\Prob[\Verifier(x, y, \pi) = 1 \mid y \neq f(x,w)\ \textnormal{ para todo }\ w \in \WitnessSpace] \leq \epsilon$, donde $\epsilon$ es un valor negligible.
		\item Eficiencia del verifier: El tiempo de ejecución de $\Verifier$ es polinomial en el tamaño de $x$ y $y$, y sublineal respecto al costo requerido para computar directamente $f(x,w)$.
	\end{enumerate}
	
	Si además el protocolo satisface la propiedad de conocimiento cero, diremos que el esquema de cómputo verificable es \ZK.
\end{definition}

En la definición anterior, asumimos que la función $f$ es fija, pública y conocida tanto por $\Prover$ como por $\Verifier$. Si la función no fuera fija, tendríamos que agregarla al conjunto de \textit{inputs} públicos, es decir que serían de la forma $(x, y, f)$.

También observamos que en la definición anterior nos referimos al esquema de cómputo verificable era \ZK, pero no hablamos de que $(\Prover, \Verifier)$ era un protocolo de pruebas ZK. Esto es porque debemos hacer una distinción fundamental antes de pasar a la siguiente sección de la tesis, entre \textit{prueba \ZK} y \textit{argumento \ZK}. Es notable cómo coloquialmente suele haber confusión entre estos dos términos y pensar que se refieren a lo mismo, usándose indistintamente. Daremos una distinción informal, basada en la distinción formal que se puede encontrar en el séptimo capítulo de \cite{Goldreich2001}.

\begin{enumerate}[noitemsep]
	\item En una prueba, la propiedad de solidez vale contra provers arbitrarios, incluso computacionalmente no acotados. Es decir, aun concediendo poder de cómputo ilimitado al prover deshonesto, la probabilidad de convencer al verifier de una afirmación falsa permanece acotada por un valor negligible.
	\item En un argumento, la propiedad de solidez sólo se exige contra provers eficientes. En particular, se cuantifica sobre provers probabilísticos en tiempo polinomial. Por eso también se habla de \textit{computationally sound proof systems}. Si además vale la propiedad de conocimiento cero, diremos que se trata de un argumento \ZK.
\end{enumerate}

Ahora bien, ¿por qué usaríamos algo más débil que una prueba \ZK, que puede ser teóricamente roto? Porque los argumentos \ZK pueden ser mucho más eficientes y escalables. En general, los argumentos \ZK asumen la imposibilidad de ciertos problemas criptográficos, como \ECDLP, \BDH, o la resistencia a colisiones de funciones de hash criptográficas.
